% ===== uphill/p2_c6.tex  (Problem 2, track 6) =====
% Written directly as a body file (not produced by tools/convert_cell.py).
% Title: Problem 2, Cell 6: $U(Q_9)$ (open question)
% Body only: packages, theorem environments and macros are in preamble.tex.
% Labels are prefixed with "p2c6:".  Uses cells 1-5 (p2c1:* ... p2c5:*).
% Bibliography (3 entries) in paper/bib/p2_c6.tex -> references.tex.

\paragraph{Official setting.} As in Cell~1 (\S\ref{p2c1:sec:count}). For C6, hand in all three
of: the value; an explicit labelling that attains it, in the same format; and a proof that no
labelling gives fewer uphill paths.

\begin{statement}
Determine $U(Q_9)$.
\end{statement}

\paragraph{Status.} \textbf{Partial.} We do \emph{not} determine $U(Q_9)$. We prove
$2328\le U(Q_9)\le2400$ (Cell~5) and reduce the remaining gap to a question about decycling sets
of $Q_9$. We conjecture $U(Q_9)=2400$ (Conjecture~\ref{p2c6:conj}), and give the evidence we have.
The labelling of Cell~5 (\S\ref{p2c5:sec:lab}) is our candidate for the optimum.

\subsubsection{What is proved}

\begin{proposition}\label{p2c6:prop:reduce}
\begin{enumerate}[label=(\alph*)]
\item $2328\le U(Q_9)\le2400$ and $227\le\nabla(Q_9)\le236$.
\item Every labelling $f$ of $Q_9$ with fewer than $2400$ uphill paths has a set $R_f$ of expensive
  vertices ($N\ge2$) that is a decycling set with at most $235$ vertices. It also satisfies
  $\sum_{u\in R_f}(N(u)-2)\,\mathrm{out}(u)\le2399-512-8|R_f|$.
\item In particular, if $\nabla(Q_9)=236$, then $U(Q_9)=2400$, attained by $L_{C_9}$.
\item Let $R$ be a decycling set of $Q_9$ contained in one parity class. Then $|R|\ge231$; and
  $|R|\ge236$ if one uses the known value $A(8,3)=20$ \cite{MS77}.
\end{enumerate}
\end{proposition}

\begin{proof}
(a) is Proposition~\ref{p2c5:prop:upper} and Theorem~\ref{p2c5:thm:main}. (b) is
Proposition~\ref{p2c1:prop:formula} with $d=9$ together with Lemma~\ref{p2c1:lem:S}(c).
(c) follows from (b) and (a). For (d), assume by symmetry that $R\subseteq E_9$, and let
$C=E_9\setminus R$. All odd vertices lie outside $R$. Two words of $C$ at distance $2$ would span a
$4$-cycle with their two (odd) common neighbours, so $C$ has pairwise distances $\ge4$. By
Lemma~\ref{p2c5:lem:del}, $|C|\le25$, so $|R|\ge256-25=231$. The maximum size of such a code equals
$A(8,3)=20$ (puncture one coordinate, respectively add a parity bit), which gives $|R|\ge236$.
\end{proof}

So a labelling beating $2400$ must come from a decycling set of $Q_9$ with at most $235$ vertices.
By (d), such a set contains vertices of \emph{both} parities, unlike every optimal set we know for
$d\le8$ (Corollary~\ref{p2c4:cor:table}).

\subsubsection{Conjecture and evidence}

\begin{conjecture}\label{p2c6:conj}
$\nabla(Q_9)=236$, and hence $U(Q_9)=2400$.
\end{conjecture}

\paragraph{Evidence (not a proof).}
\begin{itemize}
\item For $2\le d\le8$ we proved $\nabla(Q_d)=2^{d-1}-A(d-1,3)$. In each case the optimum is attained
  by a parity class minus a code (Corollary~\ref{p2c4:cor:table}). The conjecture is the same
  formula at $d=9$.
\item Up to $d=8$ this formula is squeezed between two simple lower bounds: the counting bound
  (Lemma~\ref{p2c1:lem:count}) and doubling (Lemma~\ref{p2c4:lem:double}). At $d=9$ the two bounds
  give only $225$ and $224$, while the formula gives $236$. The gap is $11$. This is why $d=9$ is the
  first open case.
\item Local search found nothing below $236$. We ran simulated annealing over sets $R$ of fixed
  size $r\in\{228,\dots,235\}$, minimising the cycle rank of $Q_9-R$. Starts were random or
  $E_9\setminus C_9$ with vertices removed. Moves swapped a vertex of $R$ with a vertex of the
  complement, and runs had up to $6\cdot10^6$ moves. No decycling set of size $\le235$ turned up; the
  best set of size $235$ left one independent cycle. A SAT search for a decycling set of size $225$
  (with the subcube bounds of Cell~5 added) did not finish within half an hour. These runs are
  exploratory and are not part of any proof.
\end{itemize}

\subsubsection{What a proof would need}

By Proposition~\ref{p2c6:prop:reduce}(c), it suffices to show that $Q_9$ has no decycling set of
size $224+s$ for $3\le s\le11$ (sizes $227$--$235$). The proof of Theorem~\ref{p2c5:thm:main}
covers $s\le2$. Two ingredients would have to be strengthened:
\begin{enumerate}[label=(\roman*)]
\item \emph{Types of light copies.} For $s\ge3$ the heavy copies of a fixed $D$ can separate
  $\{0,1\}^D$, so the light copies need not all have the same type. One would need the structure of
  decycling sets of $Q_6$ with $29$, $30$, \dots\ vertices, or copies of other dimensions, for
  instance $Q_5$ with $\nabla(Q_5)=14$ and $16$ copies per class.
\item \emph{The code bound.} Step~4 of that proof uses $|C_0|\le25$ (Lemma~\ref{p2c5:lem:del}). With
  $A(8,3)=20$ in its place, Step~4 gives $|C|\le20+\sum_{z\in Z}|C\cap N(z)|$, to be compared with
  $|C|=32-s+|Z|$. This is a contradiction while the odd part $Z$ is small. For example, when
  $|Z|\le3$ each $|C\cap N(z)|\le3$, and $|C|\le20+3|Z|<32-s+|Z|$ as soon as $2|Z|<12-s$. For a
  large $Z$ it is not.
\end{enumerate}
A decycling set of size $\le235$ would not settle the question by itself either. By
Proposition~\ref{p2c6:prop:reduce}(b), it improves on $2400$ only if the extra term
$\sum_{u\in R}(N(u)-2)\mathrm{out}(u)$ can be kept small, for instance if it is independent.

\begin{remark}[Literature]
Decycling (feedback vertex) sets of hypercubes have been studied before, e.g.\
\cite{FLP00,BB02,Pik03}. We did not use these papers and have not checked which of our values
$\nabla(Q_d)$, $d\le8$, appear there. We claim no novelty for these values. The connection with
uphill paths, the exact formula of Proposition~\ref{p2c1:prop:formula}, the resulting values of
$U(Q_d)$, and the bound $\nabla(Q_9)\ge227$ were obtained independently here.
\end{remark}
